\begin{abcliste}
\abc The electron gains the energy $e U = \frac{1}{2} m_e v^2$:
\begin{align}
 v &= \veF \\
   &= \sqrt{\frac{2\times\nce\times\U}{\ncme}} \\
   &= \ve \approx \result{\veP{0}{2}}
\end{align}

\abc
\begin{align}
 p &= \peF = \ncme\times\ve = \pe \approx \result{\peP{0}{2}}
\end{align}

\abc
\begin{align}
 \lambda &= \laF = \frac{\nch}{\pe} = \la \approx \result{\laP{-9}{2}}
\end{align}
That is about the spacing of the atoms in a crystal. (The photon formula $\lambda = hc/E$ does not hold for electrons.)

\abc $p = \sqrt{2 m_e e U}$ grows with the square root of $U$, so $\lambda = h/p \propto 1/\sqrt{U}$: four times the voltage gives twice the momentum and \result{\text{half the wavelength}}.
\end{abcliste}
